\documentclass[11pt,a4paper]{article}

\usepackage[utf8]{inputenc}
\usepackage[T1]{fontenc}
\usepackage{lmodern}
\usepackage[margin=0.76in]{geometry}
\usepackage[table]{xcolor}
\usepackage{array}
\usepackage{tabularx}
\usepackage{booktabs}
\usepackage{enumitem}
\usepackage{titlesec}
\usepackage{fancyhdr}
\usepackage{lastpage}
\usepackage{hyperref}
\usepackage[normalem]{ulem}

\definecolor{TextInk}{HTML}{233142}
\definecolor{HeadingBlue}{HTML}{1D3557}
\definecolor{AccentTeal}{HTML}{2A7F92}
\definecolor{AccentGold}{HTML}{C58B4E}
\definecolor{RuleGray}{HTML}{D6E0E5}
\definecolor{TableStripe}{HTML}{F3F8FA}
\definecolor{SoftTint}{HTML}{F7F4EE}
\definecolor{SoftBlue}{HTML}{EEF6F8}
\definecolor{SoftGreen}{HTML}{EEF7F0}
\definecolor{SoftRed}{HTML}{FBF0EF}

\hypersetup{
  colorlinks=true,
  linkcolor=HeadingBlue,
  urlcolor=AccentTeal,
  pdftitle={IB Geography SL Paper 1 Simulation Marking Guide - 2026-05-02},
  pdfauthor={IB Geography SL Preparation}
}

\color{TextInk}
\setlength{\parskip}{0.46em}
\setlength{\parindent}{0pt}
\setlength{\tabcolsep}{5pt}
\renewcommand{\arraystretch}{1.18}
\setlist[itemize]{leftmargin=1.32em,itemsep=3pt,topsep=4pt}
\setlist[enumerate]{leftmargin=1.45em,itemsep=3pt,topsep=4pt}

\newcolumntype{Y}{>{\raggedright\arraybackslash}X}
\newcolumntype{L}[1]{>{\raggedright\arraybackslash}p{#1}}

\titleformat{\section}
  {\normalfont\huge\sffamily\bfseries\color{HeadingBlue}}
  {}{0pt}{}
  [\vspace{0.12em}{\color{AccentGold}\titlerule[1pt]}]

\titleformat{\subsection}
  {\normalfont\Large\sffamily\bfseries\color{AccentTeal}}
  {}{0pt}{}

\titleformat{\subsubsection}
  {\normalfont\large\sffamily\bfseries\color{HeadingBlue}}
  {}{0pt}{}

\titlespacing*{\section}{0pt}{1.15em}{0.5em}
\titlespacing*{\subsection}{0pt}{0.85em}{0.25em}
\titlespacing*{\subsubsection}{0pt}{0.6em}{0.16em}

\pagestyle{fancy}
\setlength{\headheight}{14pt}
\fancyhf{}
\fancyhead[L]{\sffamily\color{AccentTeal}May 2 Marking Guide}
\fancyhead[R]{\sffamily\color{HeadingBlue}IB Geography SL}
\fancyfoot[C]{\sffamily\color{HeadingBlue}\thepage\ of \pageref{LastPage}}
\renewcommand{\headrulewidth}{0.4pt}

\newcommand{\term}[1]{\textcolor{HeadingBlue}{\textbf{#1}}}
\newcommand{\exam}[1]{\textcolor{AccentTeal}{\textbf{#1}}}
\newcommand{\todobox}{\fbox{\rule{0pt}{0.8em}\rule{0.8em}{0pt}}}
\newcommand{\boxnote}[3]{%
\noindent\colorbox{#1}{%
  \begin{minipage}{0.965\textwidth}
  \sffamily
  \textbf{\textcolor{HeadingBlue}{#2}} #3
  \end{minipage}
}}

\begin{document}

\begin{center}
  {\sffamily\bfseries\color{HeadingBlue}\LARGE May 2 Paper 1 Simulation Marking Guide}\\[0.25em]
  {\sffamily\color{AccentTeal}\large Self-Marking, Error Log, And May 4 Carry-Forward}
\end{center}

\boxnote{SoftRed}{Use after the 90-minute paper only.}{This guide is for
self-marking. It is not an official IB mark scheme. Award marks for accurate
geography, clear source use, valid explanation, named examples, and supported
judgment.}

\section{Strict Marking Rules}

\begin{itemize}
  \item Mark the answer that was written during the timed simulation. Do not
  improve it before scoring.
  \item For \exam{state}, give one mark for each correct, source-based point.
  \item For \exam{describe}, give credit for clear patterns, contrasts, and
  source evidence. Do not reward explanation unless it supports the description.
  \item For \exam{explain}, require a cause-and-effect chain. A listed factor
  without a link to an outcome should not receive full credit.
  \item For \exam{examine}, \exam{evaluate}, and \exam{to what extent}, require
  named evidence, more than one side, and a conclusion or judgment.
\end{itemize}

\section{Urban Environments Mark Guide}

\subsection{Question 1: State Evidence Of Uneven Environmental Quality [2]}

Award up to \term{2 marks}, one for each valid piece of evidence from Source A.
Possible answers include:

\begin{itemize}
  \item Green space is much higher in the suburban zone, at \term{34\%}, than in
  the peripheral informal settlement, at \term{3\%}.
  \item The central business district has only \term{8\%} green space and faces
  congestion and limited housing space.
  \item The peripheral informal settlement has weak service access and flood or
  pollution exposure.
  \item The inner-city regeneration zone has some green space, \term{14\%}, but
  also faces rising rents and displacement pressure.
\end{itemize}

\subsection{Question 2: Describe Two Contrasts [4]}

Award up to \term{4 marks}: up to \term{2 marks} for each clear contrast. Each
contrast should identify both zones and use Source A evidence.

\rowcolors{2}{TableStripe}{white}
\begin{tabularx}{\textwidth}{L{0.22\textwidth}Y Y}
\toprule
\textbf{Contrast type} & \textbf{Inner-city regeneration zone} & \textbf{Peripheral informal settlement} \\
\midrule
Housing and land value & High density, rising rents, reuse of former
industrial land, rent index \term{82} & Very high density, insecure housing,
rent index \term{24} \\
Environment and services & \term{14\%} green space and redevelopment pressure
& Only \term{3\%} green space, weak service access, and flood or pollution
exposure \\
Urban process & Regeneration and possible gentrification & Informal growth
and marginalization at the urban edge \\
\bottomrule
\end{tabularx}
\rowcolors{2}{}{}

\subsection{Question 3: Explain Urban Stress [4]}

Award up to \term{4 marks}: up to \term{2 marks} for each explained reason.
A full-credit reason must show the chain from urban growth to stress.

\begin{itemize}
  \item Rapid in-migration can increase housing demand faster than formal
  housing supply, causing overcrowding, informal settlement growth, or insecure
  housing.
  \item Population growth can increase pressure on water, sanitation, schools,
  transport, and healthcare, especially where local government revenue or
  planning capacity is limited.
  \item Higher densities and car dependence can increase congestion, air
  pollution, waste, noise, and urban heat-island effects.
  \item Land-value competition can displace lower-income groups from central
  areas, increasing commuting time, segregation, or social inequality.
\end{itemize}

\subsection{Question 4(a): Urban Management Strategies [10]}

Credit any valid named city. Strong examples could include Curitiba for
transport and planning, Detroit for selective regeneration and housing stress,
or another class case. Award using the level guide below.

\subsection{Question 4(b): Regeneration Or Gentrification [10]}

Credit any valid named example. Detroit is a strong anchor if the answer links
deindustrialization, population loss, selective reinvestment, gentrification,
possible displacement, and uneven benefits.

\section{Food And Health Mark Guide}

\subsection{Question 1: State Two Contrasts [2]}

Award up to \term{2 marks}, one for each valid contrast from Source B.
Possible answers include:

\begin{itemize}
  \item Region X has much higher food insecurity, at \term{38\%}, than Region Z,
  at \term{7\%}.
  \item Child undernutrition decreases from \term{28\%} in Region X to
  \term{3\%} in Region Z.
  \item Adult obesity rises from \term{6\%} in Region X to \term{32\%} in
  Region Z.
  \item Safe water access is lowest in Region X, at \term{62\%}, and highest in
  Region Z, at \term{98\%}.
\end{itemize}

\subsection{Question 2: Describe Changing Health Problems [4]}

Award up to \term{4 marks}: up to \term{2 marks} for each described pattern or
change. Good answers use Source B evidence and avoid unsupported causes.

\begin{itemize}
  \item Source B suggests a shift from undernutrition and infectious disease
  risk toward obesity and diet-related non-communicable disease risk.
  \item Region X has limited access, weak storage, high food insecurity, high
  child undernutrition, and lower safe-water access.
  \item Region Y shows a double burden: lower undernutrition than Region X but
  higher adult obesity and a mixed disease burden.
  \item Region Z has high safe-water access and low undernutrition, but high
  processed-food consumption and the highest obesity level.
\end{itemize}

\subsection{Question 3: Explain Stakeholder Effects [4]}

Award up to \term{4 marks}: up to \term{2 marks} for each explained stakeholder
chain. A full-credit point must connect the stakeholder to food security or a
health outcome.

\begin{itemize}
  \item Governments can improve food security through school meals, safety nets,
  water and sanitation, food-safety regulation, public health campaigns, or
  support for small farmers.
  \item TNCs can improve access through investment, supply chains, jobs, retail,
  and distribution, but may worsen health if marketing and pricing encourage
  high-sugar, high-salt, or ultra-processed diets.
  \item NGOs and international organizations can provide emergency food, health
  education, vaccination, clean-water projects, or disease-response support, but
  may be limited if poverty, conflict, or weak governance remain.
  \item Farmers, retailers, and consumers shape availability, price, diet, and
  waste; their choices are affected by income, land, climate risk, and policy.
\end{itemize}

\subsection{Question 4(a): TNCs In Food Security And Health [10]}

Credit any valid named TNC and place. Brazil and Nestle can be used if the
answer evaluates distribution and access benefits against processed-food
marketing, nutrition transition, and public-health concerns.

\subsection{Question 4(b): Strategies To Improve Food Security Or Health [10]}

Credit any valid named strategy or case, such as water and sanitation
improvements, school meals, disease prevention, food-security policy,
small-farmer support, food-waste reduction, or sustainable agriculture. The
answer must evaluate success and limitation.

\newpage

\section{Generic 10-Mark Level Guide}

Use this for both 10-mark answers. Award the best-fit level, then choose the
mark within the band.

\rowcolors{2}{TableStripe}{white}
\begin{tabularx}{\textwidth}{L{0.14\textwidth}Y}
\toprule
\textbf{Marks} & \textbf{Descriptor} \\
\midrule
0 & No relevant geography, blank answer, or wholly incorrect response. \\
1--2 & Basic general statements. Little or no named example. Mainly descriptive
or asserted. The command term is mostly missed. \\
3--4 & Some relevant knowledge and a partly relevant example. Explanation is
thin, uneven, or list-like. Limited balance or evaluation. \\
5--6 & Clear, mostly accurate answer with a named example and some explanation.
Shows more than one side, but evidence or judgment may be uneven. \\
7--8 & Detailed, well-structured answer with named evidence, geographic
processes, stakeholders or scale, and clear evaluation. A supported conclusion
is present. \\
9--10 & Sustained, precise, and balanced argument. Uses strong case evidence,
clear causal chains, and explicit judgment that answers the exact command term.
Limitations and uneven outcomes are handled convincingly. \\
\bottomrule
\end{tabularx}
\rowcolors{2}{}{}

\section{Command-Term Diagnostics}

\rowcolors{2}{TableStripe}{white}
\begin{tabularx}{\textwidth}{L{0.18\textwidth}Y Y}
\toprule
\textbf{Command} & \textbf{Full-credit behavior} & \textbf{Common penalty} \\
\midrule
State & Short accurate point, usually source-based & Wasting time explaining or
giving vague non-source statements \\
Describe & Pattern, contrast, trend, or feature with evidence & Explaining
causes before describing what is shown \\
Explain & Cause-and-effect chain using because / therefore logic & Listing
factors without showing how they create an outcome \\
Examine & Consider the issue from different sides before reaching a judgment &
One-sided description or no judgment \\
Evaluate & Judge success, limitation, scale, cost, equity, or sustainability &
Only describing a strategy without weighing effectiveness \\
To what extent & Make a degree judgment, not a yes/no answer & No conclusion or
conclusion that repeats the question \\
\bottomrule
\end{tabularx}
\rowcolors{2}{}{}

\section{May 2 Error Log}

Use this immediately after marking. These rows become the starting point for
the May 4 remediation block.

\begin{tabularx}{\textwidth}{L{0.17\textwidth}Y Y Y}
\toprule
\textbf{Error type} & \textbf{What I wrote or missed} & \textbf{Correct version} & \textbf{May 4 fix} \\
\midrule
Command term & \rule{0pt}{1.45cm} & & \\
Source use & \rule{0pt}{1.45cm} & & \\
Process chain & \rule{0pt}{1.45cm} & & \\
Case evidence & \rule{0pt}{1.45cm} & & \\
Evaluation & \rule{0pt}{1.45cm} & & \\
Timing & \rule{0pt}{1.45cm} & & \\
\bottomrule
\end{tabularx}

\section{Carry-Forward Decision}

\begin{tabularx}{\textwidth}{L{0.32\textwidth}Y}
\toprule
\textbf{Prompt} & \textbf{Decision} \\
\midrule
Lowest-scoring option & \rule{0.62\linewidth}{0.4pt} \\
Lowest-scoring question type & \rule{0.62\linewidth}{0.4pt} \\
One Urban content item to fix & \rule{0.62\linewidth}{0.4pt} \\
One Food and health item to fix & \rule{0.62\linewidth}{0.4pt} \\
One case-study statistic or place detail to relearn & \rule{0.62\linewidth}{0.4pt} \\
Exact practice task to redo on May 4 & \rule{0.62\linewidth}{0.4pt} \\
\bottomrule
\end{tabularx}

\vspace{0.6em}

\boxnote{SoftGreen}{May 4 priority rule.}{Redo the exact question type that
lost the most marks. If the problem was timing, redo the same option in
\term{42 minutes}. If the problem was evidence, rewrite only the 10-mark answer
with a stronger named example. If the problem was command terms, rewrite the
first sentence of every answer before adding content.}

\end{document}
